\textbf{Phần 1. Trắc nghiệm lựa chọn}
\setcounterpageref{ex}{0}

\begin{ex} %[2D4V19]
Tính kết quả của phép tính: $47 + 8$.
\choice
{$54$}
{$56$}
{\True $55$}
{$53$}
\loigiai{
$47 + 8 = 55$.
\textbf{Đáp án của câu hỏi là C}
}
\end{ex}

\begin{ex} %[2D4V19]
Lớp 2A có $28$ học sinh, hôm nay có thêm $4$ bạn học sinh mới chuyển đến. Hỏi lớp 2A có tất cả bao nhiêu học sinh?
\choice
{\True $32$ học sinh}
{$31$ học sinh}
{$33$ học sinh}
{$30$ học sinh}
\loigiai{
$28 + 4 = 32$ (học sinh).
\textbf{Đáp án của câu hỏi là A}
}
\end{ex}

\begin{ex} %[2D4V19]
Tìm $x$ biết: $x - 7 = 35$.
\choice
{$41$}
{$43$}
{$40$}
{\True $42$}
\loigiai{
$x = 35 + 7 = 42$.
\textbf{Đáp án của câu hỏi là D}
}
\end{ex}

\begin{ex} %[2D4V19]
Tính giá trị biểu thức: $36 + 9 + 5$.
\choice
{$48$}
{\True $50$}
{$49$}
{$51$}
\loigiai{
$36 + 9 = 45$; $45 + 5 = 50$.
\textbf{Đáp án của câu hỏi là B}
}
\end{ex}

\begin{ex} %[2D4V19]
Một người thợ dệt buổi sáng dệt được $38\text{ m}$ vải, buổi chiều dệt thêm $7\text{ m}$ vải. Cả ngày người thợ dệt được bao nhiêu mét vải?
\choice
{$43\text{ m}$}
{$44\text{ m}$}
{\True $45\text{ m}$}
{$46\text{ m}$}
\loigiai{
$38 + 7 = 45\text{ (m)}$.
\textbf{Đáp án của câu hỏi là C}
}
\end{ex}

\begin{ex} %[2D4V19]
Điền số thích hợp: $59 + 6 = \dots$
\choice
{\True $65$}
{$64$}
{$66$}
{$63$}
\loigiai{
$59 + 6 = 65$.
\textbf{Đáp án của câu hỏi là A}
}
\end{ex}

\begin{ex} %[2D4V19]
Cho phép cộng $\overline{a8} + 5 = 53$. Chữ số $a$ có giá trị là
\choice
{$3$}
{$5$}
{$2$}
{\True $4$}
\loigiai{
Ta có: $\overline{a8} = 53 - 5 = 48$. Do đó $a = 4$.
\textbf{Đáp án của câu hỏi là D}
}
\end{ex}

\begin{ex} %[2D4V19]
Một cây cầu dài $65\text{ m}$, người ta nối dài thêm một đoạn $8\text{ m}$. Chiều dài cây cầu sau khi nối là
\choice
{$72\text{ m}$}
{\True $73\text{ m}$}
{$74\text{ m}$}
{$71\text{ m}$}
\loigiai{
$65 + 8 = 73\text{ (m)}$.
\textbf{Đáp án của câu hỏi là B}
}
\end{ex}

\begin{ex} %[2D4V19]
Tìm tổng của số lớn nhất có hai chữ số mà chữ số hàng chục là $7$ và số lớn nhất có một chữ số.
\choice
{$86$}
{$87$}
{\True $88$}
{$89$}
\loigiai{
Số lớn nhất có 2 chữ số hàng chục là 7 là 79. Số lớn nhất có 1 chữ số là 9.\nTổng: $79 + 9 = 88$.
\textbf{Đáp án của câu hỏi là C}
}
\end{ex}

\begin{ex} %[2D4V19]
Năm nay bác Ba $56$ tuổi. Hỏi $8$ năm nữa bác Ba bao nhiêu tuổi?
\choice
{\True $64$ tuổi}
{$62$ tuổi}
{$63$ tuổi}
{$65$ tuổi}
\loigiai{
$56 + 8 = 64$ (tuổi).
\textbf{Đáp án của câu hỏi là A}
}
\end{ex}

\begin{ex} %[2D4V19]
Điền dấu thích hợp: $74 + 8 \dots 85 - 4$.
\choice
{$<$}
{$=$}
{Không so sánh được}
{\True $>$}
\loigiai{
Vế trái: $74 + 8 = 82$. Vế phải: $85 - 4 = 81$. Vì $82 > 81$ nên điền dấu $>$. 
\textbf{Đáp án của câu hỏi là D}
}
\end{ex}

\begin{ex} %[2D4V19]
Đoạn thẳng $AB$ dài $48\text{ cm}$, đoạn thẳng $BC$ dài $6\text{ cm}$. Độ dài đường gấp khúc $ABC$ là
\choice
{$52\text{ cm}$}
{\True $54\text{ cm}$}
{$53\text{ cm}$}
{$55\text{ cm}$}
\loigiai{
$48 + 6 = 54\text{ (cm)}$.
\textbf{Đáp án của câu hỏi là B}
}
\end{ex}

\textbf{Phần 2. Trắc nghiệm đúng sai}
\setcounterpageref{ex}{0}

\begin{ex} %[2D4V19]
Xét tính đúng sai của các phép cộng có nhớ số có hai chữ số với số có một chữ số:
\choiceTF
{\True Phép tính $46 + 7$ có kết quả bằng $53$.}
{Khi thực hiện phép tính $38 + 5$, chữ số hàng chục của kết quả vẫn là 3.}
{\True Biểu thức $67 + 8$ có giá trị bằng $75$.}
{Số 89 cộng với bất kì số có một chữ số nào cũng cho kết quả nhỏ hơn 95.}
\loigiai{
\begin{itemize}
    \item Ý a) Đúng: $46 + 7 = 53$.
    \item Ý b) Sai: Vì $8 + 5 = 13$ có nhớ 1 sang hàng chục nên chữ số hàng chục là $3 + 1 = 4$.
    \item Ý c) Đúng: $67 + 8 = 75$. Đúng.
    \item Ý d) Sai: Nếu cộng với 9: $89 + 9 = 98 > 95$. Khẳng định sai.
\end{itemize}
\begin{itemize}
    \item a) \textbf{Đúng}.
    \item b) \textbf{Sai}.
    \item c) \textbf{Đúng}.
    \item d) \textbf{Sai}.
\end{itemize}
}
\end{ex}

\begin{ex} %[2D4V19]
Một bác thợ rèn làm một thanh kiếm bằng cách hàn hai thanh thép: thanh dài 58 cm và thanh ngắn 7 cm:
\choiceTF
{\True Chiều dài tổng cộng của hai thanh thép là 65 cm.}
{\True Thanh dài hơn thanh ngắn 51 cm.}
{Nếu bớt ở thanh dài 8 cm thì thanh dài ngắn hơn thanh ngắn.}
{\True Tổng $58 + 7$ là một số lẻ.}
\loigiai{
\begin{itemize}
    \item Ý a) Đúng: $58 + 7 = 65\text{ cm}$.
    \item Ý b) Đúng: $58 - 7 = 51\text{ cm}$.
    \item Ý c) Sai: Thanh dài còn $58 - 8 = 50\text{ cm}$, vẫn dài hơn 7 cm.
    \item Ý d) Đúng: 65 là số lẻ. Đúng.
\end{itemize}
\begin{itemize}
    \item a) \textbf{Đúng}.
    \item b) \textbf{Đúng}.
    \item c) \textbf{Sai}.
    \item d) \textbf{Đúng}.
\end{itemize}
}
\end{ex}

\textbf{Phần 3. Trả lời ngắn}
\setcounterpageref{ex}{0}

\begin{ex} %[2D4V19]
Tính kết quả của phép tính: $68 + 7$.
\shortans{75}
\loigiai{
\textbf{Đáp số:} $75$\\
$68 + 7 = 75$.
}
\end{ex}

\begin{ex} %[2D4V19]
Một cuộn dây điện dài 45 m, người ta nối thêm 8 m dây điện nữa. Hỏi cuộn dây điện dài bao nhiêu mét?
\shortans{53}
\loigiai{
\textbf{Đáp số:} $53$\\
$45 + 8 = 53\text{ m}$.
}
\end{ex}

\begin{ex} %[2D4V19]
Tìm x biết: $x - 9 = 74$.
\shortans{83}
\loigiai{
\textbf{Đáp số:} $83$\\
$x = 74 + 9 = 83$.
}
\end{ex}

\begin{ex} %[2D4V19]
Tổng của số tròn chục liền trước 50 và số lớn nhất có một chữ số là bao nhiêu?
\shortans{49}
\loigiai{
\textbf{Đáp số:} $49$\\
Số tròn chục liền trước 50 là 40. Số lớn nhất có 1 chữ số là 9. Tổng: $40 + 9 = 49$.
}
\end{ex}

\begin{ex} %[2D4V19]
Mẹ mua 38 quả trứng gà và 7 quả trứng vịt. Hỏi mẹ mua tất cả bao nhiêu quả trứng?
\shortans{45}
\loigiai{
\textbf{Đáp số:} $45$\\
$38 + 7 = 45$ quả.
}
\end{ex}

\begin{ex} %[2D4V19]
Một hình tam giác có độ dài ba cạnh lần lượt là 25 cm, 18 cm và 7 cm. Tính chu vi hình tam giác đó.
\shortans{50}
\loigiai{
\textbf{Đáp số:} $50$\\
Chu vi: $25 + 18 + 7 = 50\text{ cm}$.
}
\end{ex}
